% ============================================================================
% CHAPTER 1: INTRODUCTION
% ============================================================================

\chapter{Εισαγωγή}
\label{ch:introduction}

% --- DRAFT: This is a starter template. Expand each section. ---

\section{Ο Κύβος του Rubik}
\label{sec:intro_cube}

Ο κύβος του Rubik, που εφευρέθηκε από τον Ούγγρο καθηγητή αρχιτεκτονικής Ernő Rubik το 1974, αποτελεί ένα από τα πιο δημοφιλή και μαθηματικά ενδιαφέροντα παζλ στην ιστορία \cite{wiki_rubiks}. Από την εμπορική του κυκλοφορία το 1980, έχουν πωληθεί εκατομμύρια αντίτυπα παγκοσμίως.

Από πλευράς επιστήμης υπολογιστών, ο κύβος παρουσιάζει εξαιρετικό ενδιαφέρον λόγω του τεράστιου χώρου καταστάσεων:

\begin{equation}
    \frac{8! \times 3^8 \times 12! \times 2^{12}}{12} = 43{,}252{,}003{,}274{,}489{,}856{,}000 \approx 4.3 \times 10^{19}
    \label{eq:state_space}
\end{equation}

Αυτός ο αριθμός υπερβαίνει κατά πολύ τον αριθμό ατόμων στο ανθρώπινο σώμα ή τα δευτερόλεπτα από τη δημιουργία του σύμπαντος, καθιστώντας την εξαντλητική αναζήτηση αδύνατη.


\section{Ορισμός του Προβλήματος}
\label{sec:intro_problem}

\subsection{Βέλτιστη Λύση}

Ως \textbf{βέλτιστη λύση} ορίζεται μια ακολουθία κινήσεων που επιλύει τον κύβο χωρίς να υπάρχει μικρότερη ακολουθία που τον επιλύει από την ίδια αρχική κατάσταση \cite{wiki_optimal}.

\subsection{God's Number}

Το 2010, οι Rokicki et al. απέδειξαν ότι κάθε κατάσταση του κύβου μπορεί να λυθεί σε το πολύ 20 κινήσεις \cite{rokicki2010gods}. Ο αριθμός αυτός είναι γνωστός ως ``God's Number'' και αποτελεί τη διάμετρο του γράφου καταστάσεων του κύβου.

\subsection{Μετρική Κινήσεων}

Στην παρούσα εργασία χρησιμοποιείται η \textbf{Half-Turn Metric (HTM)}, όπου:
\begin{itemize}
    \item Μία κίνηση = περιστροφή μιας έδρας κατά $90°$, $-90°$, ή $180°$
    \item Παραδείγματα: $R$, $R'$, $R2$ μετρούν όλα ως 1 κίνηση
\end{itemize}


\section{Στόχοι της Εργασίας}
\label{sec:intro_objectives}

Οι στόχοι της παρούσας διπλωματικής εργασίας, όπως ορίστηκαν από την εκφώνηση, είναι:

\begin{enumerate}
    \item \textbf{Υλοποίηση αλγορίθμων}: Ανάπτυξη των αλγορίθμων Thistlethwaite, Kociemba και Korf σε Python.

    \item \textbf{Εκτίμηση απόστασης}: Υλοποίηση αλγορίθμου που δέχεται μια κατάσταση και υπολογίζει πόσες κινήσεις απέχει από τη λύση.

    \item \textbf{Ευρετικές συναρτήσεις}: Εύρεση κατάλληλης ευρετικής συνάρτησης για τη βέλτιστη επίλυση με τον αλγόριθμο A* ή παραλλαγή του.
\end{enumerate}

Όλοι οι παραπάνω στόχοι επιτεύχθηκαν επιτυχώς, όπως τεκμηριώνεται στα επόμενα κεφάλαια.


\section{Συνεισφορές}
\label{sec:intro_contributions}

Οι κύριες συνεισφορές της εργασίας είναι:

\begin{itemize}
    \item Πλήρης υλοποίηση τριών αλγορίθμων επίλυσης σε Python (περίπου 8.000 γραμμές κώδικα).

    \item Ανάπτυξη \textbf{πρωτότυπης σύνθετης ευρετικής} (composite heuristic) που επιλέγει δυναμικά τη βέλτιστη στρατηγική.

    \item Πειραματική σύγκριση των αλγορίθμων με 203 test cases και benchmark data.

    \item Διαδραστικές web εφαρμογές για επίδειξη (Next.js, Streamlit).

    \item Εκπαιδευτικό υλικό σε μορφή Jupyter notebooks.
\end{itemize}


\section{Δομή της Εργασίας}
\label{sec:intro_structure}

Η εργασία οργανώνεται ως εξής:

\begin{description}
    \item[Κεφάλαιο 2] παρουσιάζει το θεωρητικό υπόβαθρο: αναπαράσταση του κύβου, στοιχεία θεωρίας ομάδων, και αλγορίθμους αναζήτησης.

    \item[Κεφάλαιο 3] περιγράφει τον αλγόριθμο Thistlethwaite (4 φάσεων).

    \item[Κεφάλαιο 4] περιγράφει τον αλγόριθμο Kociemba (2 φάσεων).

    \item[Κεφάλαιο 5] περιγράφει τον βέλτιστο αλγόριθμο Korf με pattern databases.

    \item[Κεφάλαιο 6] αναλύει την εκτίμηση απόστασης και τις ευρετικές συναρτήσεις.

    \item[Κεφάλαιο 7] παρουσιάζει την πειραματική αξιολόγηση και σύγκριση.

    \item[Κεφάλαιο 8] περιγράφει την υλοποίηση και τις web εφαρμογές.

    \item[Κεφάλαιο 9] συνοψίζει τα συμπεράσματα και προτείνει μελλοντικές επεκτάσεις.
\end{description}
